\documentclass{article}
\usepackage{amsthm}


% Language setting
% Replace `english' with e.g. `spanish' to change the document language
\usepackage[english]{babel}

% Set page size and margins
% Replace `letterpaper' with `a4paper' for UK/EU standard size
\usepackage[letterpaper,top=2cm,bottom=2cm,left=3cm,right=3cm,marginparwidth=1.75cm]{geometry}

% Useful packages
\usepackage{amsmath}
\usepackage{graphicx}
\usepackage[colorlinks=true, allcolors=blue]{hyperref}

\newtheorem{theorem}{Theorem}
\newtheorem{remark}{Remark}

\title{Cramer-Rao Lower Bound}
\author{Thanh Vuong Dang}

\begin{document}
\maketitle
Following the book \cite{Kay}, we present some versions of Cramer Rao Lower Bound theorem.
\section{Scalar version}
\begin{theorem}
Given a probability distribution function (PDF) $p(x;\theta)$ satisfy the "regularity" condition
\begin{equation}\label{1}
E\left[\dfrac{\partial ln \; p(x,\theta)}{\partial \theta}\right] = 0, \; \forall \theta,
\end{equation}
where the expectation is taken with respect to $p(x;\theta)$. Then, the variance of any unbiased estimator $\widehat{\theta}$ must not lower than 
\begin{equation}\label{2}
    \operatorname{Var}(\widehat{\theta}) \geq \dfrac{1}{-E\left[\frac{\partial^{2}ln\; p(x; \theta)}{\partial \theta^{2}}\right]},
\end{equation}
where the derivative at the true value of $\theta$ and the expectation is taken with respect to $p(x;\theta)$. Furthermore, the equality occurs for a specific unbiased estimator $\widehat{\theta} = g(x)$ if and only if
\begin{equation}
\dfrac{\partial ln\;p(x;\theta)}{\partial \theta} = I(\theta)(g(x) - \theta), \; \forall \theta,
\end{equation}
for some function $g$ and $I$. $\widehat{\theta} =  g(x)$ is called MVU estimator and the minimum variance is $\dfrac{1}{I(\theta)}.$
\end{theorem}
\begin{proof}

\end{proof}

\begin{remark}
The expectation from $\eqref{1}$ and $\eqref{2}$ is determined as
$$E\left[\dfrac{\partial ln \; p(x,\theta)}{\partial \theta}\right] = \int\dfrac{\partial ln \; p(x,\theta)}{\partial \theta}p(x;\theta)dx,$$
and 
$$E\left[\dfrac{\partial^{2} ln \; p(x,\theta)}{\partial \theta^{2}}\right] = \int\dfrac{\partial^{2} ln \; p(x,\theta)}{\partial \theta^{2}}p(x;\theta)dx.$$
\end{remark}

\section{Vector version}
\begin{theorem}
It is assumed that the PDF $p(x;\theta)$ satisfies the "regularity" conditions
\begin{equation}
E\left[\dfrac{\partial ln \; p(x,\theta)}{\partial \theta}\right] = 0, \; \forall \theta,
\end{equation}
where the expectation is taken with respect to $p(x;\theta)$. Then, the covariance matrix of any unbiased estimator $\widehat{\theta}$ satisfies 
\begin{equation}
    C_{\widehat{\theta}} - I^{-1}(\theta) \geq 0,
\end{equation}
where $\geq 0$ means that the matrix is positive semidefinite. The Fisher information matrix $$[I(\theta)]_{ij} = -E\left[\dfrac{\partial^{2}ln \; p(x;\theta)}{\partial \theta_{i}\theta_{j}}\right]$$
where the derivative are evaluated at the true value of $\theta$ and the expectation is taken with respect to $p(x;\theta)$. Furthermore, an unbiased estimator that makes $C_{\widehat{\theta}} = I^{-1}(\theta)$ if and only if
\begin{equation}
\dfrac{\partial ln \; p(x;\theta)}{\partial \theta} = I(\theta)(g(x) - \theta),
\end{equation}
for some p-dimensional function $g$ and some $p\times p$ matrix $I$. $\widehat{\theta} = g(x)$ is MVU estimator and its covariance matrix is $I^{-1}(\theta)$. 
\begin{proof}

\end{proof}
\end{theorem}
\bibliographystyle{alpha}
\bibliography{references}

\end{document}
